\section{总结与延伸}

\subsection{核心结论}

\begin{enumerate}
\item Coding 是 AGI 第二幕，因为它让模型从聊天转向执行数字世界任务。
\item 前沿模型公司的竞争窗口正在重排，上一阶段成功可能成为下一阶段路径依赖。
\item Harness Engineering 决定模型能否进入真实工作流。
\item 模型成为新一代 OS 的关键，是它能调度工具、上下文、权限和任务状态。
\item Coding/Agent 会带来白领通缩和失业窗口，社会吸收速度可能慢于技术扩散速度。
\end{enumerate}

\subsection{开放问题}

\begin{itemize}
\item Coding 模型会不会像 GPU 一样成为所有 AI 研究的基础加速器？
\item 模型 OS 会由现有大厂掌握，还是由新一代 AI-native 应用公司掌握？
\item 白领通缩会先发生在哪些行业，社会如何重新分配生产力红利？
\item 中国模型公司能否借助成本、应用和开发者生态追上第二幕？
\end{itemize}

\subsection{拓展阅读}

\begin{itemize}
\item EP139 Agent 技术史：理解 Coding/Agent 为什么是更大技术谱系的一部分。
\item EP138 罗福莉访谈：理解 Agent 范式如何改变后训练和算力分配。
\item Claude Code、Codex、OpenClaw 等工具实践：观察 Coding 模型如何改变研发工作流。
\item Harness engineering、agent evaluation、computer use agent benchmark：理解模型进入生产系统的中间层。
\end{itemize}

\begin{importantbox}{最后的判断}
如果说 2023 年的问题是“模型会不会说话”，2026 年的问题就是“模型能不能干活”。Coding 是这个转变的第一块硬地面；谁能把它接进真实工作流，谁就更接近下一代 AI 操作系统。
\end{importantbox}

\end{document}
